% Limitations and Future Work: honest gaps and the roadmap. Instruction only.

[PROCESSING INSTRUCTION (Limitations and Future Work). Be deliberately
complete and honest; the thesis depends on transparent VVUQ. Do not bury
or soften any gap. Write two short subsections.

(1) Limitations, approximations, and what could not be run. Enumerate the
gaps exactly as recorded in execution/README.md: agentic model backends
not exercised live (anthropic and openai absent, no API key, so generation
ran on deterministic offline templates); live web search not reachable, so
ingestion exhausted its 4 retries and fell back to the labeled offline
stub; PDF extraction backend (pypdf) absent, so only the size-based chunk
estimate ran; physics and robotics backends (mujoco, Isaac, ROS 2) absent
with no GPU or display; the multibyte hard-split chunking edge case that
loses 6 characters at boundaries (reported, not repaired, as it was out of
the execution scope); the LaTeX template not compiled to PDF in the
container; imagegen having no executable code in v0.1.0; and a single local
Python version (3.11.15) with cross-version 3.10 and 3.12 left to CI. Add
the two assurance gaps the Discussion raised: no independent validation
reference, and the repository's rich inputs/ corpus not yet wired into the
executed pipeline. Present these compactly in
Table~\ref{tab:limits_future}.

(2) Future work: VVUQ automation across all oncology clinical trials. Make
the central forward claim that this VVUQ automation should be adapted
through the full breadth of oncology clinical trials, because the speed,
quality, and cost benefits compound at scale while patient safety is
preserved by the gate. Lay out the tractable next steps in priority order:
supply an independent validation reference; wire the inputs/ corpus through
ingestion so codegen and papergen run on real source material; make the
triple simulation genuinely stochastic with more runs and confidence
intervals; enable the live model, network, PDF, and physics paths; repair
the multibyte chunker and add a multibyte test; close all three thesis legs
(codegen, imagegen, papergen) end to end; and emit machine readable
results. Frame the destination as a flexible, future proof assurance layer
that lets the industry check these steps off once rather than repeatedly.

SOURCE DATA to render as a publication table (paired now and next):

\begin{table}[H]
\caption{Limitations now and the tractable next steps.}
\label{tab:limits_future}
\centering
\begin{tabular}{L{6.4cm} L{7.0cm}}
\toprule
\textbf{Limitation (this run)} & \textbf{Future step} \\
\midrule
  No independent validation reference (self referential accept) & Supply an external gold standard so validation carries weight. \\
  inputs/ corpus not wired into the pipeline & Feed the chunked corpus through ingestion into codegen and papergen. \\
  Generation on offline templates; no live backends & Enable live model, web search, PDF, and physics paths. \\
  Near-deterministic 3 run simulation (CV 0.0007) & Genuine stochasticity, more runs, confidence intervals, sensitivity. \\
  Multibyte hard-split loses 6 characters & Carry trailing bytes forward; add a multibyte chunking test. \\
  imagegen has no executable code in v0.1.0; template uncompiled & Close codegen, imagegen, and papergen legs end to end. \\
  Single local Python version & Run the 3.10 / 3.11 / 3.12 matrix locally, emit JSON results. \\
\bottomrule
\end{tabular}
\end{table}

SOURCE FILES: execution/README.md (the ``Limitations, approximations, and
what could not be run'' section) is the authoritative list;
execution/04-stage1-automation/README.md for the multibyte detail.

HOW TO PROCESS: keep the limitations factual and specific, reference
Table~\ref{tab:limits_future} rather than repeating it, and let the future
work read as a confident roadmap toward all-trial adoption. Cite
\cite{repo-cancer-automated}, \cite{sel2025vvuq}, \cite{kawchak2026federated},
\cite{kawchak2025qsppdac}, and \cite{fda2026realtime}.]
